\documentclass{article}
\usepackage{amsmath,amssymb,amsthm}
\newtheorem{theoreme}{Théorème}
\newtheorem{definition}{Définition}
\newtheorem{proposition}{Proposition}
\begin{document}

\section{Éléments propres}

Dans tout le chapitre, $E$ est un $\mathbb{K}$-espace vectoriel de dimension finie $n \geq 1$
et $u \in \mathcal{L}(E)$. On note $A \in \mathcal{M}_n(\mathbb{K})$ une matrice.

\begin{definition}
Un scalaire $\lambda \in \mathbb{K}$ est une valeur propre de $u$ s'il existe $x \neq 0$ tel que
$u(x) = \lambda x$. Un tel $x$ est un vecteur propre associé à $\lambda$.
Le sous-espace propre associé est $E_\lambda(u) = \ker(u - \lambda \mathrm{id}_E)$.
L'ensemble des valeurs propres est le spectre, noté $\mathrm{Sp}(u)$.
\end{definition}

\begin{proposition}
Des sous-espaces propres associés à des valeurs propres deux à deux distinctes sont en somme
directe. En particulier, une famille de vecteurs propres associés à des valeurs propres
distinctes est libre.
\end{proposition}

\begin{proof}
Par récurrence sur le nombre de valeurs propres, en appliquant $u - \lambda_p \mathrm{id}$
à une relation de dépendance.
\end{proof}

\section{Polynôme caractéristique}

\begin{definition}
Le polynôme caractéristique de $A$ est $\chi_A(X) = \det(X I_n - A)$. Il est unitaire de
degré $n$, et $\chi_A(X) = X^n - \mathrm{tr}(A) X^{n-1} + \dots + (-1)^n \det(A)$.
Deux matrices semblables ont le même polynôme caractéristique.
\end{definition}

\begin{proposition}
$\lambda$ est valeur propre de $A$ si et seulement si $\chi_A(\lambda) = 0$.
Si $\lambda$ est racine de $\chi_A$ de multiplicité $m(\lambda)$, alors
$1 \leq \dim E_\lambda \leq m(\lambda)$.
\end{proposition}

\section{Diagonalisation}

\begin{definition}
$u$ est diagonalisable s'il existe une base de $E$ formée de vecteurs propres de $u$.
De façon équivalente, $A$ est diagonalisable s'il existe $P$ inversible et $D$ diagonale
telles que $A = P D P^{-1}$.
\end{definition}

\begin{theoreme}[Critère de diagonalisabilité]
Les assertions suivantes sont équivalentes :
\begin{enumerate}
  \item $u$ est diagonalisable ;
  \item $E = \bigoplus_{\lambda \in \mathrm{Sp}(u)} E_\lambda(u)$ ;
  \item $\sum_{\lambda \in \mathrm{Sp}(u)} \dim E_\lambda(u) = n$ ;
  \item $\chi_u$ est scindé sur $\mathbb{K}$ et pour toute valeur propre $\lambda$,
        $\dim E_\lambda(u) = m(\lambda)$.
\end{enumerate}
\end{theoreme}

\begin{proposition}
Si $\chi_u$ est scindé à racines simples (en particulier si $u$ a $n$ valeurs propres
distinctes), alors $u$ est diagonalisable. La réciproque est fausse : $\mathrm{id}_E$.
\end{proposition}

Méthode pratique : calculer $\chi_A$, en trouver les racines, puis pour chaque valeur propre
résoudre $(A - \lambda I_n) X = 0$ pour obtenir une base de $E_\lambda$ ; $P$ a pour colonnes
ces vecteurs.

\section{Polynômes annulateurs}

\begin{theoreme}[Cayley-Hamilton]
Le polynôme caractéristique annule $u$ : $\chi_u(u) = 0$.
\end{theoreme}

\begin{definition}
Le polynôme minimal $\pi_u$ est l'unique polynôme unitaire qui engendre l'idéal des
polynômes annulateurs de $u$. Il divise $\chi_u$ et a les mêmes racines.
\end{definition}

\begin{theoreme}
$u$ est diagonalisable si et seulement si $u$ admet un polynôme annulateur scindé à racines
simples, si et seulement si $\pi_u$ est scindé à racines simples.
\end{theoreme}

\section{Trigonalisation}

\begin{theoreme}
$u$ est trigonalisable si et seulement si $\chi_u$ est scindé sur $\mathbb{K}$.
En particulier, toute matrice de $\mathcal{M}_n(\mathbb{C})$ est trigonalisable.
\end{theoreme}

Si $A$ est trigonalisable de valeurs propres $\lambda_1, \dots, \lambda_n$ (comptées avec
multiplicité), alors $\mathrm{tr}(A) = \sum_i \lambda_i$ et $\det(A) = \prod_i \lambda_i$.

\section{Réduction des matrices symétriques réelles}

\begin{theoreme}[Théorème spectral]
Toute matrice symétrique réelle $S \in \mathcal{S}_n(\mathbb{R})$ est diagonalisable dans une
base orthonormée : il existe $P \in O_n(\mathbb{R})$ et $D$ diagonale réelle telles que
$S = P D P^{\top}$.
\end{theoreme}

\end{document}
